\documentclass[11pt]{article}
\newcommand{\courseshort}{Data Structures (Make-up) · 010153103}
\newcommand{\lecturetitle}{L10: Graphs}
\usepackage{lecture_style}

\begin{document}
\lectureheader{Lecture 10: Graphs --- BFS, DFS \& Shortest Path}{Pre-class brief}{Goodrich Ch.~14}

\begin{objbox}
\begin{itemize}[leftmargin=*]
  \item Define a graph $G=(V,E)$ and distinguish directed/undirected, weighted/unweighted.
  \item Represent a graph with an adjacency list and an adjacency matrix.
  \item Apply BFS (queue) and DFS (stack/recursion) to traverse a graph.
  \item Trace Dijkstra's algorithm on a small weighted graph.
\end{itemize}
\end{objbox}

\section*{1. Graphs}

A graph $G = (V, E)$ has $n = |V|$ vertices and $m = |E|$ edges. Edges may be \emph{directed} (one-way) or \emph{undirected} (two-way), and may carry \emph{weights}.

\textbf{Adjacency list} (best for sparse graphs): for each vertex $v$, store a list of its neighbors. Memory $O(n+m)$.

\textbf{Adjacency matrix} (best for dense graphs): an $n\times n$ matrix $A$ with $A[u][v]=1$ if edge $(u,v)\in E$. Memory $O(n^2)$.

\section*{2. BFS --- shortest paths in unweighted graphs}

Use a queue and a \texttt{visited} set. Visit vertices in order of distance from the source.

\begin{lstlisting}
void bfs(List<List<Integer>> adj, int s) {
    boolean[] seen = new boolean[adj.size()];
    Queue<Integer> q = new ArrayDeque<>();
    seen[s] = true; q.add(s);
    while (!q.isEmpty()) {
        int u = q.remove();
        visit(u);
        for (int v : adj.get(u)) {
            if (!seen[v]) { seen[v] = true; q.add(v); }
        }
    }
}
\end{lstlisting}

\textbf{Cost:} $O(n + m)$.

\section*{3. DFS --- explore deeply}

Use the call stack (recursion). Used for connectivity, cycle detection, topological sort.

\begin{lstlisting}
void dfs(List<List<Integer>> adj, int u, boolean[] seen) {
    seen[u] = true;
    visit(u);
    for (int v : adj.get(u)) if (!seen[v]) dfs(adj, v, seen);
}
\end{lstlisting}

\textbf{Cost:} $O(n + m)$.

\section*{4. Dijkstra (single-source shortest path, non-negative weights)}

Maintain $\texttt{dist}[v]$, initially $\infty$ except $\texttt{dist}[s]=0$. Use a min-priority queue keyed by tentative distance.

\begin{enumerate}[leftmargin=*]
  \item Pull the closest unfinalized vertex $u$.
  \item For every neighbor $v$, \emph{relax}: if $\texttt{dist}[u] + w(u,v) < \texttt{dist}[v]$, update $\texttt{dist}[v]$.
  \item Repeat until all vertices finalized.
\end{enumerate}

\textbf{Cost:} $O((n+m)\log n)$ with a heap-based PQ.

\begin{exbox}{Pre-Class Exercises (do BEFORE the lecture)}
\textbf{Exercise 10.1 --- Two representations.} For the undirected graph with edges
\verb|{(A,B), (A,C), (B,D), (C,D), (D,E)}|, draw the adjacency list and write the $5\times 5$ adjacency matrix. Order vertices A--E.

\textbf{Exercise 10.2 --- BFS trace.} Run BFS from \texttt{A} on the graph above. List the visit order and the BFS-tree distances $\text{dist}(A, \cdot)$.

\textbf{Exercise 10.3 --- DFS trace.} Run DFS from \texttt{A}. Visit neighbors in alphabetical order. List the visit order.

\textbf{Exercise 10.4 --- Dijkstra by hand.} On the weighted directed graph
\[
A\xrightarrow{1}B,\quad A\xrightarrow{4}C,\quad B\xrightarrow{2}C,\quad B\xrightarrow{5}D,\quad C\xrightarrow{1}D,\quad D\xrightarrow{3}E
\]
run Dijkstra from \texttt{A}. Show the \texttt{dist} table after each pull from the PQ.

\textbf{Exercise 10.5 --- Why no negative edges?} Construct a tiny graph (3 vertices) with one negative edge where Dijkstra produces a wrong shortest path.

\textbf{Exercise 10.6 --- Application.} Given a 2D grid with walls, formulate the problem ``shortest number of moves from $S$ to $G$'' as a graph problem. Which traversal would you use, and why?
\end{exbox}

\begin{disbox}
\begin{itemize}[leftmargin=*]
  \item When does adjacency list beat adjacency matrix in practice? When does the matrix win?
  \item Why does BFS find shortest paths in unweighted graphs but not in weighted ones?
  \item Discuss Exercise 10.5 --- this is exactly why algorithms like Bellman--Ford exist.
  \item Wrap-up: how does this whole course --- arrays, lists, queues, heaps, hash maps --- power Dijkstra?
\end{itemize}
\end{disbox}

\vspace{1em}
\noindent
{\color{primary}\rule{\linewidth}{0.4pt}}\\
\textbf{End of make-up course.} Practice problems: re-implement each data structure from scratch and time it on $n = 10^3, 10^5, 10^7$.

\end{document}
